% Avsnitt 13.1-13.6, ur lectures/L13/appendix/a_crc15.md.

\section{Gränssnitt}
\label{c13:sec:iface}

\modulefig[0.62]{lectures/L13/appendix/images/crc15.png}

\begin{booktable}{@{}p{5.6em}p{3.4em}p{5.2em}L@{}}
  \thead{Port} & \thead{Riktn.} & \thead{Typ} & \thead{Betydelse} \\
  \midrule
  \code{clock} & in & \code{std_logic} & 50 MHz systemklocka. \\
  \code{reset_s2_n} & in & \code{std_logic} & Aktiv låg, synkroniserad reset. \\
  \code{clear} & in & \code{std_logic} & Återför registret till idel nollor vid den här flanken,
    utan en reset. \\
  \code{enable} & in & \code{std_logic} & Integrera \code{data} i registret vid den här flanken när
    \code{'1'}. \\
  \code{data} & in & \code{std_logic} & Den aktuella biten, biten som sänds (TX) eller nyss samplats
    av bussen (RX). \\
  \code{crc} & out & \code{crc_t} & Registrets aktuella 15-bitars CRC-värde. \\
  \code{valid} & out & \code{std_logic} & \code{'1'} närhelst registret för närvarande läser idel
    nollor. \\
\end{booktable}

Inga generics. \code{crc15_tb} binder sina portar \textbf{positionellt}, så ordningen och typerna
ovan är det som måste stämma; namnen är dina. \code{crc_t} (en 15-bitars vektor), dess bredd
\code{CRC_WIDTH} och polynomet \code{CRC_POLY} kommer från \code{can_def} (\chapref{c10:ch}).

% ------------------------------------------------------------------------------------------
\section{Rekursionen}
\label{c13:sec:recursion}

Håll ett 15-bitars register (\code{crc_t}), nollställt till idel nollor av \code{reset_s2_n}. Vid
varje klockflank där \code{enable = '1'}, integrera en bit:

\begin{textdrawing}
feedback = crc(CRC_WIDTH-1) xor data      -- Inkommande bit XOR registrets översta bit.
shifted  = crc skiftat ett steg vänster, ny lägsta bit '0'
crc      = shifted xor CRC_POLY           -- om feedback = '1'
         = shifted                        -- annars
\end{textdrawing}

Det är hela motorn, femton vippor och en skiftning XOR-grindad av polynomets avgreningar. När
\code{enable} är \code{'0'}, håll registret oförändrat. Driv \code{crc} kontinuerligt från registret,
och \code{valid} som ”registret är idel nollor”.

\begin{warning}
\code{clear} gör med registret vad \code{reset_s2_n} gör, men den är ett \textbf{synkront kommando,
inte en reset}: den hör hemma inuti den klockade grenen, prövad före \code{enable} så att den vinner
när båda är höga vid samma flank. Vik inte in den i det asynkrona resetvillkoret; det skulle göra
\code{clear} asynkron och nivåkänslig, en annan konstruktion som råkar passera testbänken (som alltid
håller \code{clear} över en hel klockcykel). Gränssnittstabellen säger ”vid den här flanken” och
menar det.
\end{warning}

\code{clear} finns eftersom motorn återanvänds för varje ram och håller inget eget tillstånd per ram:
\code{can_controller} pulsar den en gång i början av varje ram (\chapref{c16:ch}). En ram som får
löpa till slut återför registret till noll av sig själv, genom
generera-sedan-kontrollera-egenskapen nedan, så på den lyckliga vägen ändrar \code{clear} ingenting.
En ram som \emph{avbryts} halvvägs gör inte det, och det är det fallet den finns till för.

% ------------------------------------------------------------------------------------------
\section{En motor, två jobb}
\label{c13:sec:twojobs}

\begin{itemize}
  \item \textbf{Generera:} aktivera motorn över SOF till och med datafältets slut. Läs sedan
    \code{crc}; det är värdet som ska sändas som ramens CRC-fält.
  \item \textbf{Kontrollera:} aktivera den över SOF till och med datafältet \emph{och fortsätt
    sedan}, och mata in det \emph{mottagna} CRC-fältets egna bitar också. Om ramen anlände intakt
    driver det registret tillbaka till exakt noll, en egenskap hos aritmetiken, inte något motorn
    blivit tillsagd att förvänta sig. \code{valid} rapporterar precis det.
    \code{can_controller} (\chapref{c17:ch}) kontrollerar \code{valid} bara vid just den punkten; den
    behöver ingen andra, separat kontrollkrets.
\end{itemize}

Båda jobben förbrukar \textbf{bara riktiga bitar}. Stoppbitar (\chapref{c14:ch}) kommer aldrig in i
CRC:n på någondera sidan av länken; anroparens grindning av \code{enable} är det som håller dem ute
(\chapref{c16:ch} och \ref{c17:ch} bygger precis det), och motorn själv vet aldrig att de finns. Det
är därför \code{crc15} inte behöver någon kunskap om ramfält eller stoppning: den som driver
\code{enable} har redan tagit de besluten.

Båda jobben förutsätter att registret börjar en ram på noll, vilket är exakt vad \code{clear}
garanterar. Lägg märke till hur mycket som vilar på det: en nod som överger en ram halvvägs (den
förlorar arbitreringen, eller ser en stoppningsöverträdelse, eller underkänns i CRC-kontrollen, allt
sådant som \chapref{c17:ch} gör) slutar mata motorn med ett \emph{partiellt} värde i den. Utan en
nollställning per ram får nästa ram den resten som utgångsvärde, så noden beräknar en CRC ingen annan
håller med om och beräknar en felaktig kontroll på allt den tar emot, för varje ram från då till
reset. En enda avbruten ram skulle ta noden ur bussen permanent.

\code{valid} läser \code{'1'} vid reset också, före någon data alls; meningslöst där, meningsfullt
bara vid kontrollpunkten ovan.

% ------------------------------------------------------------------------------------------
\section{Ett genomräknat exempel, för hand}
\label{c13:sec:worked}

Börja med det enklaste fallet. Mata in en enda \code{1} i en nyss resettad motor: \code{feedback} är
\code{0 xor 1 = 1}, \code{shifted} är fortfarande idel nollor, så \code{crc} blir
\code{0 xor CRC_POLY}, polynomet självt. Det är det första testbänken kontrollerar, och en bra
rimlighetskontroll på din egen kod.

Nu fyra bitar, \code{1001}, MSB först. \code{top} är registrets bit 14 \emph{före} flanken, och
\code{CRC_POLY} är \code{100010110011001}:

\begin{textdrawing}
         top  data  fb   shifted          xor poly   crc
reset      -     -   -   -                       -   000000000000000
bit 1      0     1   1   000000000000000        ja   100010110011001
bit 2      1     0   1   000101100110010        ja   100111010101011
bit 3      1     0   1   001110101010110        ja   101100011001111
bit 4      1     1   0   011000110011110       nej   011000110011110
\end{textdrawing}

Läs en rad i taget så slutar motorn se ut som aritmetik. Varje flank gör samma tre saker: XOR:a den
inkommande biten med biten som faller av toppen, skifta vänster, och XOR:a tillbaka polynomet bara om
det resultatet var \code{1}.

Två rader är värda att stanna vid. \textbf{Bit 2} matar in en \code{0} och utlöser ändå XOR:en,
eftersom \code{1}:an som lämnar toppen är det som sätter \code{feedback}; den inkommande biten är
bara halva beslutet. \textbf{Bit 4} matar in en \code{1} och utlöser den \emph{inte}, eftersom den
översta biten också är \code{1} och de två tar ut varandra. Det är därför en CRC svarar på \emph{var}
en bit sitter och inte bara på hur många ettor det finns, och det är den egenskap
Övning~\ref{c13:ex:missed} använder för att fånga en förvanskning som en enkel summa inte kan se.

Lägg också märke till att varje rads \code{shifted}-kolumn är föregående rads \code{crc} skiftad ett
steg vänster, med en \code{0} som kommer in nedifrån. Ingenting adderas någonsin i den låga änden;
registrets innehåll kommer helt och hållet från polynomet som XOR:as in om och om igen vid olika
förskjutningar.

% ------------------------------------------------------------------------------------------
\section{Vad testbänken låser fast}
\label{c13:sec:tb}

\begin{itemize}
  \item Strax efter reset är \code{crc} idel nollor och \code{valid = '1'}.
  \item En enda \code{'1'}-bit ger \code{crc = CRC_POLY}.
  \item \code{clear} återför ett laddat register till idel nollor utan att \code{reset_s2_n} rörs.
  \item En 19-bitars sekvens av SOF, ID, RTR och kontrollfält ger en CRC som stämmer med ett värde
    beräknat oberoende i Python (inte härlett ur den här VHDL-koden).
  \item Att mata den nyss beräknade CRC:ns egna bitar rakt in igen återför registret till idel nollor
    (\code{valid = '1'}), generera-sedan-kontrollera-egenskapen.
\end{itemize}

% ------------------------------------------------------------------------------------------
\section{Att lägga in den i \code{can_controller}}
\label{c13:sec:wire}

Samma tvådelade ändring som \chapref{c12:ch}:s, ett block längre in. I
\file{controller/can_controller.vhd}, lägg till signalerna under \code{bit_timer}-gruppen:

\begin{vhdlcode}
    -- crc15.
    signal crc_clear, crc_enable, crc_data_in, crc_valid : std_logic;
    signal crc_value                                     : crc_t;
\end{vhdlcode}

och instansen under \code{bit_timer1}:

\begin{vhdlcode}
    crc1: entity work.crc15
        port map(clock, reset_s2_n, crc_clear, crc_enable, crc_data_in, crc_value, crc_valid);
\end{vhdlcode}

Lägg märke till att \code{crc15} kopplas till ingenting annat än klockan, resetten och sina egna
signaler: den tar sina bitar från vilket skiftregister som än är aktivt, och båda de registren är
ännu oskrivna. Kopplingen som förenar dem är \chapref{c16:ch}:s, och den är ett enda uttryck snarare
än en port map.
